\chapter{Conclusioni}
\label{chap:conclusion}
L'esperienza di tirocinio svolta ha permesso di completare con successo l'intero ciclo di vita del software per la modernizzazione dei servizi di latenza. Il passaggio dall'architettura PHP/Python al nuovo ecosistema in Go ha portato notevoli benefici prestazionali e strutturali.
\begin{itemize}
    \item \emph{Efficienza Computazionale}: Il passaggio a binari nativi compilati in contenitori Debian-Slim ha ridotto notevolmente l'utilizzo della memoria RAM rispetto al sistema basato su Apache, PHP e Python.
    \item \emph{Documentazione}: la scrittura della documentazione in parallelo con lo sviluppo del codice ha permesso una dettagliata descrizione di tutte le componenti del sistema che facilitano eventuali modifiche, fix e aggiunte future.
    \item \emph{Robustezza operativa}: La gestione manuale delle chiamate ping garantiscono una alta affidabilità rispetto al sistema preesistente.
    \item \emph{Trade-off velocità/affidabilità}: L'aggiunta di diversi layer aggiuntivi rispetto all'esecuzione nativa del ping dal sistema operativo riporta dei leggeri rallentamenti (circa $50-70 \mu s$).
\end{itemize}

Questo rallentamento è però trascurabile, poiché la latenza richiede una precisione delle misurazioni nell'ordine dei millisecondi. Inoltre il rallentamento è visibile solo nel caso specifico dell'esecuzione del ping sull'interfaccia di loop-back, una richiesta possibile ma che ha poco senso ai fini dello strumento.

\section{Sviluppi Futuri}vediamo tutta questa perdita sugli hop intermedi, ma non su
Sebbene l'obiettivo principale del progetto fosse la migrazione del servizio esistente verso la nuova infrastruttura, nel corso dello sviluppo sono emerse alcune opportunità di estensione che meritano di essere documentate in vista di sviluppi futuri. In particolare, è stata individuata la possibilità di introdurre una nuova funzionalità di diagnostica di rete, descritta di seguito.
\subsubsection{Traceroute}
Una delle feature individuate come possibile estensione futura riguarda l'implementazione di un meccanismo di \emph{traceroute} applicativo, pensato per tracciare il percorso di rete attraversato dalle richieste dell'utente.

L'idea alla base consiste nel lanciare, dai nodi Anycast globali, un traceroute di rete che esponga informazioni sui passaggi (hop) dall'\ac{IP} sorgente, in modo analogo a quanto avviene con il comando \texttt{traceroute} a livello di sistema operativo.
Questo permetterebbe di ottenere una visione più granulare della topologia applicativa, utile sia in fase di debug che di monitoraggio delle latenze introdotte dai singoli nodi attraversati.

Un aspetto rilevante, emerso proprio grazie al lavoro di migrazione svolto, è che il sistema attuale è stato progettato secondo criteri di modularità ed estendibilità.
Questo si traduce in un vantaggio concreto per l'implementazione della nuova feature: l'introduzione del traceroute sul server di downstream non richiederebbe una riprogettazione del sistema, ma si limiterebbe all'aggiunta di una nuova classe che sfrutta il wrapper \ac{ICMP} già progettato e integrato nell'architettura esistente. In altre parole, l'infrastruttura di base necessaria, ovvero la componente in grado di gestire e interpretare i pacchetti \ac{ICMP}, è già disponibile, e il lavoro richiesto si ridurrebbe principalmente all'esposizione di questa funzionalità tramite la nuova route, senza impattare in modo significativo sul resto del sistema.

Questa caratteristica conferma la bontà delle scelte architetturali adottate durante la fase di migrazione, che hanno privilegiato un disaccoppiamento chiaro tra i vari livelli del servizio, rendendo agevole l'aggiunta di nuove funzionalità anche in un secondo momento, senza dover intervenire pesantemente sul codice preesistente.